\chapter{Desarrollo Empírico y Resultados}
\label{ch:resultados}

Este capítulo contrasta empíricamente las hipótesis planteadas en el Capítulo \ref{ch:introduccion} mediante econometría de series de tiempo. El paso desde el análisis descriptivo del Capítulo \ref{ch:descriptivo} hacia la estimación econométrica permite evaluar los determinantes de la Función de Reacción Fiscal argentina sobre la ventana muestral única (1996--2025, $n=120$, treinta años exactos): DF-GLS, Johansen, VECM, DOLS, IV-2SLS, Hansen, Bai-Perron, el SVAR restringido y el Threshold VECM se estiman todos sobre esta misma ventana (Secciones \ref{sec:dfgls}--\ref{sec:tvecm_resultados}), con el VECM, el SVAR restringido y el Threshold VECM como técnica de referencia para la Función de Reacción Fiscal y el protocolo de no linealidad, en reemplazo de DOLS, a pedido del director de tesis (Secciones \ref{sec:vecm} y \ref{sec:svar_metodologia} del Capítulo \ref{ch:metodologia}). La evidencia recogida es mayormente débil o estadísticamente no concluyente: la regla de reacción fiscal lineal de Bohn no opera de forma estable bajo ninguna de las técnicas aplicadas. Dicha inoperancia estadística constituye el hallazgo central del capítulo y se analiza en profundidad en el Capítulo \ref{ch:discusion}, junto con la evidencia de un régimen crítico de riesgo soberano que el modelo de umbral de Hansen no confirma como fatiga fiscal en su especificación en niveles, sino como el patrón inverso (Sección \ref{sec:hansen}). La descomposición de varianza del SVAR restringido, por su parte, cambia de manera sustantiva respecto de lo que sostenía una estimación previa sobre una ventana más corta (Sección \ref{sec:svar_resultados}), y esa reversión se discute con el mismo criterio de transparencia que el resto del capítulo.

\section{Contrastes Adicionales de Raíz Unitaria (DF-GLS)}
\label{sec:dfgls}

La Sección \ref{sec:estacionariedad} del Capítulo \ref{ch:descriptivo} ya avanzó un diagnóstico de integración basado en los contrastes complementarios de ADF y KPSS. Para robustecerlo se aplica adicionalmente el test de raíz unitaria \textit{Dickey-Fuller Generalized Least Squares} (DF-GLS) de \textcite{elliott1996}, que remueve la tendencia determinística mediante mínimos cuadrados generalizados antes de estimar la regresión aumentada, logrando una potencia estadística sustancialmente mayor que el ADF convencional en muestras acotadas como la presente ($n=120$).

\begin{table}[H]
\centering
\caption{\label{tab:dfgls} Resultados del Test de Raíz Unitaria DF-GLS (Elliott-Rothenberg-Stock, 1996), Ventana Única (1996--2025).}
\resizebox{\textwidth}{!}{%
\begin{tabular}{lrrrrl}
\toprule
\textbf{Variable} & \textbf{Estad. (nivel)} & \textbf{Crít. 5\% (nivel)} & \textbf{Estad. (1.\ dif.)} & \textbf{Crít. 5\% (1.\ dif.)} & \textbf{Conclusión} \\
\midrule
$pb_t$ & -1.30 & -3.01 & -5.98 & -2.12 & I(1) \\
$d_t$ & -2.65 & -3.01 & -4.07 & -2.12 & I(1) \\
EMBI+ ($risk_t$) & -2.30 & -3.01 & -9.16 & -2.11 & I(1) \\
$\tilde{y}_t$ & -3.46 & -3.02 & -4.93 & -2.12 & I(0) \\
\bottomrule
\end{tabular}%
}
\begin{flushleft}
\footnotesize \textit{Nota.} Elaboración propia (\texttt{codigo/modelos/fase1\_estacionariedad.py}). Test DF-GLS especificado con constante y tendencia determinística en niveles, y únicamente constante en primeras diferencias.
\end{flushleft}
\end{table}

El DF-GLS confirma, con mayor potencia estadística que el protocolo del Capítulo \ref{ch:descriptivo}, el carácter integrado de orden uno de $pb_t$, $d_t$ y del riesgo soberano: las tres no rechazan la raíz unitaria en niveles y sí la rechazan con holgura tras la primera diferencia. Para la deuda en particular, DF-GLS es más informativo que el ADF simple de la Sección \ref{sec:estacionariedad}: mientras el ADF por sí solo dejaba una lectura ambigua entre I(0) e I(1) para $d_t$ en esta ventana (Sección \ref{sec:ventana_ampliada} del Capítulo \ref{ch:descriptivo}), el DF-GLS, de mayor potencia, no rechaza la raíz unitaria en niveles ($-2.65$ frente a un crítico de $-3.01$) y respalda con más autoridad el tratamiento I(1) que el resto del protocolo mantiene por continuidad. La brecha del producto, en cambio, rechaza la raíz unitaria tanto en niveles como en diferencias, I(0) sin ambigüedad. Esta clasificación, tres series I(1) y una I(0), es la base sobre la que la Sección \ref{sec:johansen_cap6} retoma el contraste formal de cointegración de Johansen para evaluar si el sistema conjunto admite una relación de equilibrio de largo plazo.

\section{Relación de Largo Plazo: Cointegración de Johansen}
\label{sec:johansen_cap6}

Como se estableció en la Sección \ref{sec:dfgls}, el sistema conformado por la ratio Deuda/PIB, el Resultado Primario, el EMBI+ y el TCRM es I(1). Esta circunstancia obliga a verificar formalmente, mediante el procedimiento de máxima verosimilitud de Johansen, si existe al menos una combinación lineal estacionaria entre estas cuatro variables, condición necesaria para que una regresión en niveles no sea espuria.

\begin{table}[H]
\centering
\caption{\label{tab:johansen} Test de cointegración de Johansen (Traza y Máximo Autovalor), sistema [Deuda/PIB, Result. Primario, EMBI+, TCRM], Ventana Única (1996--2025).}
\begin{tabular}{lrrcrrc}
\toprule
\multirow{2}{*}{$H_0$: rango $\le$} & \multicolumn{3}{c}{\textbf{Estadístico de Traza}} & \multicolumn{3}{c}{\textbf{Estadístico de Máx. Autovalor}} \\
 & \textbf{Valor} & \textbf{Crít. 95\%} & \textbf{Rechaza} & \textbf{Valor} & \textbf{Crít. 95\%} & \textbf{Rechaza} \\
\midrule
0 & 45.18 & 47.85 & No & 19.97 & 27.59 & No \\
1 & 25.21 & 29.80 & No & 15.10 & 21.13 & No \\
2 & 10.10 & 15.49 & No & 7.03  & 14.26 & No \\
3 & 3.07  & 3.84  & No & 3.07  & 3.84  & No \\
\bottomrule
\end{tabular}
\begin{flushleft}
\footnotesize \textit{Nota.} Elaboración propia (\texttt{codigo/modelos/fase23\_reestimacion\_1996\_2025.py}), $k\_ar\_diff=1$ (BIC). La secuencia formal de decisión se detiene en la primera hipótesis no rechazada; al no rechazarse $H_0: r=0$, la evaluación se detiene ahí y la significatividad reportada en $r\ge 1$ se muestra solo por completitud de la tabla.
\end{flushleft}
\end{table}

Sobre la ventana única de treinta años, ambos estadísticos de Johansen coinciden en no rechazar $H_0: r=0$: la Traza ($45.18$) queda por debajo de su crítico al 95\% ($47.85$) y el Máximo Autovalor ($19.97$) también queda por debajo del suyo ($27.59$). Es un diagnóstico más contundente que el de una ventana más corta, aunque en el sentido opuesto al que ofrecía una estimación previa sobre 2004--2025: allí ambos estadísticos coincidían en rechazar $r=0$ y encontrar exactamente una relación de cointegración; aquí coinciden en no encontrar ninguna. Esta tensión, ya documentada en la Sección \ref{sec:resultados_vecm} y en la Sección \ref{sec:etapa2_johansen} del Capítulo \ref{ch:metodologia}, no se oculta: se impone igualmente un rango $r=1$ por motivo teórico, la relación de largo plazo de \citet{bohn1998}, para poder estimar el VECM y el SVAR restringido (Secciones \ref{sec:resultados_vecm} y \ref{sec:svar_resultados}), documentando explícitamente que la imposición no es una elección neutral.

El pedido explícito del director de tesis de estimar un VAR restringido o VECM en lugar de DOLS (Sección \ref{sec:vecm} del Capítulo \ref{ch:metodologia}) sostiene la estimación de un Vector de Corrección de Errores como técnica de referencia de esta tesis. La Sección \ref{sec:resultados_vecm} reporta dicha estimación, junto con la selección del orden de rezagos del VAR subyacente y la sensibilidad del rango de cointegración a esa elección; DOLS se conserva en las Secciones \ref{sec:dols_resultados}--\ref{sec:sensibilidad_temporal} como ejercicio de robustez de ecuación única sobre la misma ventana, coherente con el argumento de \citet{bohn2007} de que la verificación empírica de la condición de transversalidad de Bohn admite más de un camino de estimación válido.

\section{Estimación de la Función de Reacción Fiscal por VECM (Ventana Ampliada)}
\label{sec:resultados_vecm}

Sobre la ventana ampliada 1996T1--2025T4 ($n=120$, treinta años exactos, construida mediante el empalme de la Sección \ref{sec:empalme} del Capítulo \ref{ch:metodologia}), se estima un Modelo de Vectores con Corrección de Error (VECM) para el mismo sistema de cuatro variables [Deuda/PIB, Resultado Primario, EMBI+, TCRM], con la brecha del producto como variable exógena estacionaria. A diferencia de DOLS, que trata a $pb_t$ como variable dependiente y purga la endogeneidad de corto plazo mediante adelantos y rezagos \textit{ad hoc} de $\Delta d_t$, el VECM modela el sistema completo como conjuntamente endógeno: la endogeneidad queda integrada directamente en la estructura del modelo, sin necesidad de una construcción auxiliar.

\subsection{Selección de Rezagos y Sensibilidad del Rango de Cointegración}

La selección de rezagos del VAR en niveles sobre la ventana ampliada no arroja unanimidad entre criterios: AIC y FPE seleccionan 5 rezagos, HQIC selecciona 2, y BIC, el más parsimonioso, selecciona 1. Dado que $n=120$ es una muestra moderada para un sistema de cuatro variables (cada rezago adicional cuesta 16 parámetros), se adopta el rezago del criterio BIC ($k\_ar\_diff=1$) como especificación de referencia, reportando en la Tabla \ref{tab:vecm_sensibilidad_lags} la sensibilidad de la estimación a esta elección.

Más consecuente aún es la sensibilidad del propio contraste de cointegración al punto de inicio de la muestra, reportada en la Tabla \ref{tab:vecm_sensibilidad_inicio}: con el rezago del criterio BIC, el estadístico de Traza no rechaza la ausencia de cointegración ($H_0: r=0$) para ningún inicio de muestra anterior a 2000T1, y sí la rechaza a partir de esa fecha.

\begin{table}[H]
\centering
\caption{\label{tab:vecm_sensibilidad_inicio} Sensibilidad del Contraste de Traza de Johansen al Inicio de la Muestra (Ventana Ampliada, $k\_ar\_diff=1$).}
\begin{tabular}{lrrrc}
\toprule
\textbf{Inicio de la muestra} & \textbf{$n$} & \textbf{Traza ($H_0: r=0$)} & \textbf{Crítico 95\%} & \textbf{Rango hallado} \\
\midrule
1996T1 (completa) & 120 & 45.18 & 47.85 & 0 \\
1997T1 & 116 & 45.46 & 47.85 & 0 \\
1998T1 & 112 & 46.24 & 47.85 & 0 \\
1999T1 & 108 & 46.38 & 47.85 & 0 \\
2000T1 & 104 & 50.95 & 47.85 & 1 \\
2001T1 & 100 & 58.62 & 47.85 & 1 \\
\bottomrule
\end{tabular}
\begin{flushleft}
\footnotesize \textit{Nota.} Elaboración propia (\texttt{codigo/modelos/fase23\_reestimacion\_1996\_2025.py}). Los valores desde 1999T1 en adelante coinciden exactamente con los ya reportados para la ventana ampliada original de 108 observaciones, verificando la continuidad metodológica de la extensión.
\end{flushleft}
\end{table}

A diferencia de la ventana ampliada de 108 observaciones, donde excluir un único año revertía el diagnóstico, en la ventana de treinta años el resultado es más contundente: ningún punto de inicio entre 1996T1 y 1999T1, cuatro alternativas distintas, rechaza la ausencia de cointegración; recién al excluir también 1998 y 1999 completos (inicio en 2000T1) el estadístico de Traza supera al crítico. Esta fragilidad es coherente con la literatura sobre la distorsión del contraste de Johansen ante quiebres estructurales severos no modelados \citep{gregory1996,johansen2000}, y es precisamente el tipo de tensión que el pedido del director de extender la ventana ``tan atrás como fuera posible'' buscaba hacer emerger: treinta años de muestra no ofrecen, por sí solos, una imagen más nítida de la cointegración de largo plazo del sistema, sino una imagen más consistentemente ambigua a lo largo de un rango más amplio de puntos de inicio posibles. Se impone en consecuencia un rango de cointegración $r=1$ por motivo teórico, el mismo criterio ya aplicado en la Sección \ref{sec:johansen_cap6}, documentando explícitamente que la imposición no es una elección neutral.

\subsection{Estimación VECM y Sensibilidad a la Elección de Rezagos}

\begin{table}[H]
\centering
\caption{\label{tab:vecm_beta_alpha} Estimación VECM, Ventana Ampliada ($n=120$, $k\_ar\_diff=1$, rango de cointegración $=1$).}
\begin{tabular}{lrrrr}
\toprule
\multicolumn{5}{l}{\textit{Panel A. Vector de cointegración $\boldsymbol{\beta}$ (normalizado sobre deuda\_pib)}} \\
\midrule
\textbf{Variable} & \textbf{Coef.} & \textbf{Error Est.} & \textbf{$z$} & \textbf{$p$-valor} \\
\midrule
deuda\_pib & 1.0000 & --- & --- & --- \\
pb\_pib & 5.9653 & 2.595 & 2.299 & 0.022 \\
EMBI & -0.0150 & 0.003 & -4.930 & 0.000 \\
TCRM & 2.0466 & 10.160 & 0.201 & 0.840 \\
const & -45.3547 & 16.944 & -2.677 & 0.007 \\
\midrule
\multicolumn{5}{l}{\textit{Panel B. Coeficientes de ajuste $\boldsymbol{\alpha}$ (velocidad de corrección hacia el equilibrio)}} \\
\midrule
\textbf{Ecuación} & \textbf{Coef.} & \textbf{Error Est.} & \textbf{$z$} & \textbf{$p$-valor} \\
\midrule
deuda\_pib & -0.0736 & 0.026 & -2.854 & 0.004 \\
pb\_pib & -0.0000 & 0.003 & -0.005 & 0.996 \\
EMBI & 6.3560 & 2.837 & 2.240 & 0.025 \\
TCRM & -0.0003 & 0.001 & -0.416 & 0.677 \\
\bottomrule
\end{tabular}
\begin{flushleft}
\footnotesize \textit{Nota.} Elaboración propia (\texttt{codigo/modelos/fase23\_reestimacion\_1996\_2025.py}). Deterministic="ci" (constante restringida al espacio de cointegración). $\beta$.1 (deuda\_pib) fijo en 1 por normalización, sin error estándar.
\end{flushleft}
\end{table}

El coeficiente de ajuste de interés directo, $\alpha_{pb}$, la velocidad a la que el resultado primario corrige un desvío del período anterior respecto de la relación de largo plazo, es esencialmente nulo ($-0.00002$) y no significativo ($p=0.996$): no hay evidencia, bajo esta especificación, de que el resultado primario reaccione sistemáticamente para restablecer el equilibrio de largo plazo del sistema, un resultado incluso más categórico que el de la ventana ampliada de 108 observaciones ($\alpha_{pb}=0.0044$, $p=0.204$). El coeficiente de ajuste de la propia deuda es, nuevamente, significativo y con el signo correcto ($\alpha_{deuda}=-0.0736$, $p=0.004$): es la deuda la que corrige hacia el equilibrio de largo plazo. A diferencia de la ventana de 108 observaciones, aquí el riesgo soberano también emerge como significativo ($\alpha_{EMBI}=6.356$, $p=0.025$): ante un desvío del sistema respecto de su relación de largo plazo, es el EMBI+ el que absorbe buena parte del ajuste, no el resultado primario. Leído junto al Capítulo \ref{ch:discusion}, este resultado refuerza, con evidencia adicional y sobre una muestra un 11\% mayor, la lectura de esta tesis: el ajuste del sistema está dominado por el canal de riesgo soberano y por revaluaciones de stock, no por el esfuerzo fiscal corriente.

\begin{table}[H]
\centering
\caption{\label{tab:vecm_sensibilidad_lags} Sensibilidad de la Estimación VECM a la Elección de Rezagos (Ventana Ampliada, rango de cointegración $=1$).}
\begin{tabular}{lrrrl}
\toprule
\textbf{$k\_ar\_diff$} & \textbf{Rezago VAR} & \textbf{$\beta_{deuda}$ (normalizado s/ $pb$)} & \textbf{$\alpha_{pb}$} & \textbf{$p$-valor ($\alpha_{pb}$)} \\
\midrule
1 (BIC) & 2 & 0.1677 & -0.00002 & 0.996 \\
2 & 3 & 0.0614 & -0.0047 & 0.022$^{*}$ \\
3 & 4 & 0.0949 & -0.0034 & 0.288 \\
4 & 5 & 0.0870 & -0.0019 & 0.460 \\
5 (AIC, FPE) & 6 & 0.1025 & +0.0019 & 0.558 \\
\bottomrule
\end{tabular}
\begin{flushleft}
\footnotesize \textit{Nota.} $^{*}p<0.05$. Elaboración propia (\texttt{codigo/modelos/fase23\_reestimacion\_1996\_2025.py}).
\end{flushleft}
\end{table}

La Tabla \ref{tab:vecm_sensibilidad_lags} muestra, en la ventana ampliada de treinta años, un patrón parcialmente distinto al de la ventana de 108 observaciones: el signo de $\beta_{deuda}$ ahora sí es estable, positivo en las cinco especificaciones, a diferencia de la inestabilidad de signo que exhibía antes. La significatividad de $\alpha_{pb}$, en cambio, sigue sin ser robusta: es significativa en una única especificación aislada ($k\_ar\_diff=2$, $p=0.022$), pero con signo negativo, lo que implicaría que el resultado primario se aleja del equilibrio en lugar de corregirlo hacia él, un signo económicamente perverso para una regla de reacción estabilizadora. Ninguna otra especificación, incluida la de referencia por BIC, acompaña ese hallazgo. Sigue siendo, como antes, el patrón típico de un hallazgo aislado por azar de especificación antes que de un efecto robusto: la lectura que se sostiene en las cinco especificaciones es la ausencia de evidencia consistente de una reacción fiscal estabilizadora vía VECM, el mismo diagnóstico que DOLS entrega sobre la ventana original (Sección \ref{sec:dols_resultados}).

\subsection{Robustez: Velocidad de Ajuste por Administración Presidencial}
\label{sec:dummies_presidenciales}

A pedido del director de tesis, se evaluó si la velocidad de ajuste $\alpha_{pb}$ del VECM difiere sistemáticamente por administración presidencial, como contraste complementario al umbral de riesgo soberano de Hansen (Sección \ref{sec:hansen}). El ejercicio inicial interactuó variables \textit{dummy} de gobierno con la deuda rezagada en niveles ($d_{t-1}$) mediante MCO; se descartó esa especificación porque $d_{t-1}$ es I(1) (Sección \ref{sec:dfgls}), y una interacción con una variable no estacionaria arriesga confundir el efecto de administración con la propia tendencia de la deuda, precisamente el problema de regresión espuria que el resto de este capítulo evita mediante DOLS y VECM. En su lugar, se interactúan las \textit{dummy} de administración con el término de corrección de error del VECM de referencia, $\widehat{ect}_{t-1} = \hat{\boldsymbol{\beta}}' \mathbf{x}_{t-1}$ (Tabla \ref{tab:vecm_beta_alpha}), que es estacionario por construcción al ser la combinación lineal cointegrante, y se estima $\Delta pb_t = c + \alpha_{pb}\, \widehat{ect}_{t-1} + \sum_i \delta_i\, (D_i \cdot \widehat{ect}_{t-1}) + \gamma\, \tilde{y}_t + \varepsilon_t$ por MCO con errores HAC, sobre la ventana de 88 observaciones (2004--2025, series reales; pendiente de reestimarse sobre la ventana única), con Néstor Kirchner (2004T1--2007T3) como categoría de referencia y las cinco administraciones siguientes (CFK I, CFK II, Macri, Alberto Fernández, Milei) como \textit{dummy}.

\begin{table}[H]
\centering
\caption{\label{tab:dummies_presidenciales} Velocidad de Ajuste ($\alpha_{pb}$) Implícita por Administración Presidencial.}
\begin{tabular}{lrr}
\toprule
\textbf{Administración} & \textbf{$\alpha_{pb}$ implícito} & \textbf{$p$-valor (interacción)} \\
\midrule
Néstor Kirchner (ref., 2004T1--2007T3) & -0.0002 & 0.867 \\
Cristina Fernández I (2007T4--2011T3) & 0.0080 & 0.128 \\
Cristina Fernández II (2011T4--2015T3) & 0.0044 & 0.440 \\
Macri (2015T4--2019T3) & 0.0037 & 0.410 \\
Alberto Fernández (2019T4--2023T3) & 0.0240 & 0.001$^{***}$ \\
Milei (2023T4--2025T4) & 0.0174 & 0.115 \\
\bottomrule
\end{tabular}
\begin{flushleft}
\footnotesize \textit{Nota.} $\alpha_{pb}$ implícito = coeficiente de $\widehat{ect}_{t-1}$ más el coeficiente de la interacción administración$\times \widehat{ect}_{t-1}$ correspondiente. El $p$-valor reportado es el de la interacción individual, no el del nivel. Test F conjunto de las cinco interacciones: $F=5.402$, $p=0.0003$. Errores HAC (Newey-West, 4 rezagos). $^{***}p<0.01$. Elaboración propia (\texttt{codigo/modelos/fase18\_dummies\_presidenciales.py}).
\end{flushleft}
\end{table}

El test F conjunto rechaza la hipótesis de homogeneidad entre administraciones ($p=0.0003$), pero individualmente solo una es significativa: Alberto Fernández ($\alpha_{pb}=0.024$, $p=0.001$), con una velocidad de ajuste positiva y estadísticamente sólida hacia el equilibrio de largo plazo del sistema. Las cinco administraciones restantes, incluida la de referencia, no son individualmente distinguibles de cero. Este resultado no debe leerse como evidencia de una regla de reacción fiscal activa y estable durante ese gobierno en particular (el diagnóstico de fondo del capítulo, sostenido por DOLS y por el VECM en ambas ventanas, sigue siendo la ausencia de una reacción sistemática de muestra completa), sino como una fuente adicional, y acotada a un único período, de la inestabilidad paramétrica ya documentada en la Sección \ref{sec:sensibilidad_temporal}: el rechazo del test conjunto es compatible con que la velocidad de ajuste no sea constante en el tiempo, sin que ello implique que la mayoría de los gobiernos individuales exhiban, por sí solos, una reacción fiscal estadísticamente identificable.

\section{Estimación de la Función de Reacción Fiscal por Mínimos Cuadrados Ordinarios Dinámicos (DOLS)}
\label{sec:dols_resultados}

Ante el diagnóstico de a lo sumo una relación de cointegración ($r=1$) y la inestabilidad del contraste de Traza frente al orden de rezagos, documentados en la Sección \ref{sec:johansen_cap6}, se opta por Mínimos Cuadrados Ordinarios Dinámicos (DOLS, \citealp{stock1993}) por sobre su alternativa semi-paramétrica, \textit{Fully Modified OLS} (FMOLS, \citealp{phillips1990}): ambos estimadores son asintóticamente equivalentes y corrigen el sesgo de segundo orden de la endogeneidad de corto plazo, pero DOLS incorpora dicha corrección \textit{ex ante}, mediante adelantos y rezagos de las primeras diferencias de los regresores, mientras que FMOLS la aplica \textit{ex post} sobre los residuos, requiriendo estimar no paramétricamente un espectro de frecuencia cero mediante un núcleo y un ancho de banda arbitrarios que se vuelven inestables en muestras acotadas. \citet{stock1993} documentan mediante simulación de Monte Carlo que esta diferencia se traduce en mejores propiedades de tamaño y potencia para DOLS precisamente cuando el número de observaciones es limitado (el caso de esta tesis, $n=120$, que se reduce a $n=110$ tras incorporar la ampliación dinámica de adelantos y rezagos), lo que motiva la elección por una ventaja de eficiencia en muestras finitas.

\subsection{Alcance de la Ecuación de Largo Plazo: por qué el DOLS no replica el sistema de Johansen}
\label{sec:alcance_dols}

El contraste de la Sección \ref{sec:johansen_cap6} evalúa el espacio de cointegración de un sistema de cuatro variables (Deuda/PIB, Resultado Primario, EMBI+ y TCRM), mientras que la ecuación DOLS que sigue retiene únicamente $d_{t-1}$ y la brecha del producto $\tilde{y}_t$ como regresores estructurales. Esta asimetría es deliberada y responde a tres razones que conviene explicitar.

En primer lugar, ambos procedimientos responden a preguntas distintas. El contraste de Johansen es un diagnóstico \textit{de sistema}: verifica si existe alguna combinación lineal estacionaria entre las cuatro variables, condición necesaria para que una regresión en niveles no sea espuria. La ecuación DOLS, en cambio, estima un \textit{objeto teórico específico y previamente definido}: la Función de Reacción Fiscal de \citet{bohn1998}, cuya forma canónica es $pb_t = f(d_{t-1}, \tilde{y}_t)$. El parámetro de interés $\rho$ está definido dentro de esa especificación; ampliarla con regresores adicionales estimaría un objeto distinto y no comparable con la literatura de referencia \citep{bohn1998, ghosh2013}.

En segundo lugar, y de manera decisiva, el EMBI+ y el TCRM no son exógenos respecto del resultado primario, sino \textit{simultáneamente determinados} con él. El test de Wu-Hausman de la Sección \ref{sec:iv2sls} rechaza la exogeneidad del riesgo soberano de manera categórica ($p<0.001$). Incorporar el EMBI+ como regresor adicional en la ecuación de largo plazo, por tanto, no añadiría un control: introduciría un regresor endógeno cuya presencia sesgaría el propio coeficiente $\rho$ que la ecuación busca identificar. Por esa razón el riesgo soberano no se traslada a la ecuación DOLS sino a un diseño de estimación separado, el IV-2SLS de la Sección \ref{sec:iv2sls}, que es el procedimiento apropiado para un regresor endógeno. El TCRM, por su parte, opera en esta tesis como determinante del stock de deuda por la vía del Ajuste Stock-Flujo (Sección \ref{sec:restriccion_intertemporal}), es decir, sobre la variable explicativa $d_{t-1}$, y no como determinante directo del esfuerzo primario contemporáneo; su inclusión en la ecuación de reacción confundiría el canal de valuación con el canal de comportamiento fiscal.

En tercer lugar, el diagnóstico de multicolinealidad respalda que esta restricción no descarta información relevante a costa de sesgo por variable omitida grave. La especificación ampliada (Bohn + EMBI+ + TCRM) arroja un VIF máximo de $1.91$ y un número de condición de $2.38$ sobre regresores estandarizados, ambos holgadamente por debajo de los umbrales convencionales de $10$ y $30$: los regresores del núcleo de Bohn ($VIF=1.02$) son prácticamente ortogonales entre sí y comparten poca varianza con las dos variables excluidas. La estimación restringida no está, en consecuencia, absorbiendo por vía de la constante un efecto sustantivo de las variables omitidas.

Cabe precisar, además, que todas las variables del sistema ($pb_t$, $d_t$, $\tilde{y}_t$ y $risk_t$) se emplean en niveles (o, en el caso de $\tilde{y}_t$, en desviación porcentual de su tendencia), y no en logaritmos. Esta decisión no es arbitraria: el Resultado Primario ($pb_t$) y la Brecha del Producto ($\tilde{y}_t$) toman valores negativos en una proporción sustancial de la muestra, déficits fiscales y fases recesivas del ciclo respectivamente, lo cual vuelve matemáticamente inadmisible su transformación logarítmica. Mantener todas las variables en niveles, en lugar de aplicar logaritmos únicamente a los regresores que lo permitirían (como $d_t$ o $risk_t$), preserva la interpretación homogénea y directa de los coeficientes estimados como puntos porcentuales del PIB (o puntos básicos) por unidad de la variable explicativa, y evita el problema de interpretación mixta que resultaría de combinar elasticidades (variables en logaritmos) con efectos marginales en niveles dentro de una misma ecuación estructural.

La Tabla \ref{tab:dols} reporta la estimación DOLS de la Función de Reacción Fiscal, especificada como $pb_t = \alpha + \rho\, d_{t-1} + \gamma\, \tilde{y}_t + \sum_{j=-4}^{4} \phi_j \Delta d_{t-j} + \varepsilon_t$, donde los términos $\Delta d_{t-j}$ (adelantos y rezagos de la primera diferencia de la deuda) constituyen la ampliación dinámica de Stock-Watson que garantiza la súper-consistencia del estimador de $\rho$ frente a la endogeneidad de corto plazo, sin ser en sí mismos de interés estructural directo.

\begin{table}[H]
\centering
\caption{\label{tab:dols} Estimación DOLS de la Función de Reacción Fiscal, Ventana Única ($n=110$, errores HAC, 5 rezagos).}
\begin{tabular}{lrrrr}
\toprule
\textbf{Variable} & \textbf{Coef.} & \textbf{Error Est. (HAC)} & \textbf{Estadístico z} & \textbf{$p$-valor} \\
\midrule
Constante & -1.3325 & 0.850 & -1.567 & 0.117 \\
$d_{t-1}$ (ratio deuda rezagada) & 0.0159 & 0.013 & 1.232 & 0.218 \\
Brecha del Producto ($\tilde{y}_t$) & 0.0271 & 0.040 & 0.680 & 0.496 \\
$\Delta d_t$ & 0.0183 & 0.035 & 0.524 & 0.601 \\
$\Delta d_{t-1}$ & -0.0176 & 0.024 & -0.731 & 0.465 \\
$\Delta d_{t+1}$ & -0.0338 & 0.036 & -0.950 & 0.342 \\
$\Delta d_{t-2}$ & 0.0086 & 0.021 & 0.407 & 0.684 \\
$\Delta d_{t+2}$ & 0.0174 & 0.024 & 0.734 & 0.463 \\
$\Delta d_{t-3}$ & 0.0005 & 0.022 & 0.022 & 0.982 \\
$\Delta d_{t+3}$ & -0.0203 & 0.025 & -0.820 & 0.412 \\
$\Delta d_{t-4}$ & 0.0320 & 0.042 & 0.755 & 0.450 \\
$\Delta d_{t+4}$ & -0.0124 & 0.035 & -0.356 & 0.721 \\
\midrule
\bottomrule
\end{tabular}
\begin{flushleft}
\footnotesize \textit{Nota.} $R^2=0.134$. $R^2$ ajustado$=0.037$. $F=0.780$ ($p=0.659$). Durbin-Watson$=0.117$. Errores estándar robustos a heterocedasticidad y autocorrelación (HAC, Newey-West, 5 rezagos). Elaboración propia (\texttt{codigo/modelos/fase3\_reaccion\_fiscal.py}).
\end{flushleft}
\end{table}

El coeficiente de interés estructural, $\rho$, que mide la respuesta de largo plazo del superávit primario ante la deuda heredada, resulta de $0.0159$ y no es estadísticamente significativo ($p=0.218$). La muestra disponible no permite rechazar la hipótesis nula de ausencia de reacción fiscal sistemática. A diferencia de la estimación sobre la ventana de 88 observaciones, aquí la estimación puntual presenta signo positivo, el requerido por la condición de sostenibilidad de \citet{bohn1998}: la reversión de signo no cambia el veredicto de fondo (el coeficiente sigue siendo estadísticamente nulo), pero sí matiza la lectura de un signo ``incorrecto'' que sostenía la ventana más corta. Ninguno de los términos de ampliación dinámica alcanza significatividad individual en esta ventana, a diferencia de la ventana de 88 observaciones, donde $\Delta d_{t+1}$ sí era significativo al 1\%. Se reitera que estos coeficientes no admiten una interpretación estructural directa, son parámetros no estructurales (\textit{nuisance parameters}) de la técnica DOLS.

El estadístico Durbin-Watson de 0.117, muy por debajo del valor de referencia de ausencia de autocorrelación (2.0) y más bajo aún que el 0.321 de la ventana de 88 observaciones, indica autocorrelación residual positiva de primer orden más pronunciada, atribuible a la inercia del proceso presupuestario argentino: una porción sustancial del gasto público (previsional, salarial, transferencias indexadas) y de la recaudación exhibe una dinámica autorregresiva que una especificación DOLS de rezagos acotados no absorbe por completo. La prueba de multiplicador de Lagrange de Breusch-Godfrey (cuatro rezagos) confirma esta persistencia, rechazando de manera categórica la hipótesis nula de no autocorrelación ($p<10^{-19}$). Los errores estándar HAC (Newey-West) empleados en la Tabla \ref{tab:dols} corrigen la inferencia frente a esta propiedad, de modo que la significatividad reportada para $\rho$ y los términos de ampliación dinámica permanece válida pese a la persistencia residual.

Como complemento del diagnóstico de autocorrelación serial, se aplicó adicionalmente el test ARCH-LM de \citet{engle1982} para detectar heterocedasticidad condicional en los residuos del DOLS, es decir, conglomerados de volatilidad (\textit{volatility clustering}), un patrón habitual en series fiscales y financieras de economías con historia de crisis recurrentes. El estadístico de multiplicador de Lagrange, con cuatro rezagos, resulta en $LM=66.29$ ($p<10^{-12}$), rechazando categóricamente la hipótesis nula de homocedasticidad condicional, con mayor contundencia aún que en la ventana de 88 observaciones ($LM=18.74$, $p=0.0009$). Este hallazgo refuerza la decisión metodológica de reportar la totalidad de las estimaciones de este capítulo bajo errores estándar robustos a heterocedasticidad y autocorrelación (HAC, Newey-West): la combinación de autocorrelación serial (Durbin-Watson, Breusch-Godfrey) y heterocedasticidad condicional (ARCH-LM) documentada en esta sección indica que cualquier inferencia basada en errores estándar convencionales habría subestimado sistemáticamente la incertidumbre de los coeficientes estimados.

Los diagnósticos de estabilidad estructural CUSUM/CUSUMSQ y la reestimación DOLS por subperíodo, reportados más abajo en esta sección y en la Sección \ref{sec:sensibilidad_temporal} para la ventana de 88 observaciones, no se recalcularon todavía sobre la ventana única de treinta años: se presentan tal como fueron corridos, identificados explícitamente como pendientes de esa actualización, en lugar de reportar cifras no verificadas.

Finalmente, se evaluó la estabilidad estructural de la especificación DOLS mediante el test de Suma Acumulada de Residuos Recursivos (CUSUM) y su variante cuadrática (CUSUMSQ) de \citet{brown1975}, que testean si los parámetros del modelo permanecen constantes a lo largo de la muestra sin necesidad de especificar \textit{a priori} una fecha de quiebre. La Figura \ref{fig:cusum} presenta ambos estadísticos junto a sus bandas de confianza al 95\%.

\begin{figure}[H]
\centering
\caption{Test de Estabilidad Estructural CUSUM y CUSUMSQ de la Especificación DOLS}
\label{fig:cusum}
\includegraphics[width=0.95\textwidth]{figura_6_2_cusum.png}
\par\vspace{2mm}
{\footnotesize\textit{Nota.} Residuos OLS recursivos sobre la misma especificación de la Tabla \ref{tab:dols}. Bandas de confianza al 95\% según \citet{brown1975} (CUSUM) y \citet{harvey1990} (CUSUMSQ). Elaboración propia.}
\end{figure}

Ambos estadísticos permanecen dentro de sus respectivas bandas de confianza a lo largo de la totalidad de la muestra (desviación máxima del CUSUMSQ de 0.393, frente a un límite de 0.596 al 5\%), sin cruzar el límite en ningún punto. Este resultado admite una lectura que exige matización cuidadosa: la ausencia de inestabilidad detectada por el CUSUM/CUSUMSQ no contradice la evidencia de quiebres estructurales documentada en la Sección \ref{sec:quiebres_estructurales} del Capítulo \ref{ch:descriptivo} ni la motivación teórica del modelo de umbral de Hansen (Sección \ref{sec:hansen}), sino que testea una propiedad distinta: mientras el CUSUM evalúa la constancia de los \textit{parámetros estimados} de una especificación lineal única a lo largo del tiempo, los quiebres detectados por \texttt{ruptures} y el propio umbral de Hansen refieren a un cambio de \textit{régimen} en el nivel de las series o en la relación condicional entre ellas, una hipótesis alternativa distinta de la constancia paramétrica de la regresión lineal DOLS, y no necesariamente en tensión con ella. En otras palabras: los coeficientes de la Tabla \ref{tab:dols}, aunque débiles y no significativos en la especificación OLS que el CUSUM evalúa, no cruzan los límites de confianza en ningún punto de la muestra: el no-rechazo del CUSUM implica ausencia de evidencia de inestabilidad en la especificación lineal OLS, no una afirmación positiva de estabilidad del coeficiente estructural $\rho$. La Sección \ref{sec:sensibilidad_temporal} complementa este diagnóstico mostrando que, bajo ampliación dinámica DOLS por subperíodo, la reacción fiscal estimada es positiva y significativa en ambas mitades de la muestra (2004--2014: $\rho=0.0373$, $p=0.0015$; 2015--2025: $\rho=0.0530$, $p=0.0026$, ambos bajo la especificación estática; véase la Tabla \ref{tab:sensibilidad_temporal} para las versiones DOLS), siendo el coeficiente nulo de muestra completa un artefacto de la ampliación global que absorbe correlaciones de corto plazo de distinta naturaleza en cada régimen.

Este patrón de autocorrelación serial persistente en los residuos de la regla de Bohn confirma dos realidades metodológicas de las economías periféricas: en primer lugar, las perturbaciones sobre el resultado primario (como caídas en la recaudación por comercio exterior o rigideces estructurales del gasto corriente) poseen una memoria estocástica de largo plazo; en segundo lugar, la persistencia residual indica que una especificación lineal estática resulta incapaz de absorber la totalidad de la varianza del vector de solvencia, actuando como un fuerte indicador indirecto de la presencia de no linealidades o quiebres estructurales ocultos que justifican el salto analítico hacia la regresión de umbral de Hansen implementada en la Sección \ref{sec:hansen}.

\section{Corrección por Endogeneidad: Variables Instrumentales en Dos Etapas (IV-2SLS)}
\label{sec:iv2sls}

El EMBI+ no puede tratarse como una variable de control strictly exógena: la teoría económica predice causalidad inversa, dado que un deterioro del resultado primario eleva la prima de riesgo percibida por los acreedores, mientras que un incremento exógeno del EMBI+ encarece el financiamiento y presiona al resultado fiscal. Para atender este problema, se instrumentó el EMBI+ mediante el índice de volatilidad global VIX y un diferencial de riesgo soberano regional (EMBI$_{Brasil}$, variable de aproximación de mercado construida sobre el ETF EMB -iShares JPMorgan USD Emerging Markets Bond, que replica el JPMorgan EMBI Global Diversified-), en reemplazo del Tipo de Cambio Real Multilateral (TCRM) rezagado empleado en una versión preliminar de este ejercicio, descartado por violar la restricción de exclusión (impacta las cuentas fiscales de manera directa e independiente del canal de riesgo soberano, vía retenciones a las exportaciones y subsidios indexados al tipo de cambio real). Un diferencial de riesgo soberano regional captura el mismo canal de riesgo sistémico de mercados emergentes que castiga al EMBI+ argentino, sin ese efecto presupuestario directo sobre $pb_t$.

\begin{table}[H]
\centering
\caption{\label{tab:iv2sls} Estimación IV-2SLS de la Función de Reacción Fiscal, Ventana Única ($n=119$).}
\resizebox{\textwidth}{!}{%
\begin{tabular}{lrrrr}
\toprule
\textbf{Variable} & \textbf{Coef.} & \textbf{Error Est. (HAC)} & \textbf{$t$} & \textbf{$p$-valor} \\
\midrule
Constante & -0.0513 & 0.713 & -0.072 & 0.943 \\
$d_{t-1}$ & -0.0313 & 0.024 & -1.317 & 0.188 \\
Brecha del Producto & 0.0578 & 0.034 & 1.705 & 0.088$^{*}$ \\
EMBI+ (instrumentado) & 0.0012 & 0.001 & 2.259 & 0.024$^{**}$ \\
\midrule
\multicolumn{5}{l}{\textit{Estadísticos de diagnóstico instrumental}} \\
\midrule
F (1.\textsuperscript{a} etapa, relevancia) & \multicolumn{4}{l}{$15.87$ \quad ($p<0.001$) \quad [Umbral Staiger-Stock: $>10$] $\checkmark$} \\
Test Wu-Hausman (endogeneidad) & \multicolumn{4}{l}{$2.17$ \quad ($p=0.144$) \quad [No rechaza exogeneidad] $\times$} \\
Test Sargan (sobreidentificación) & \multicolumn{4}{l}{$0.006$ \quad ($p=0.939$) \quad [Instrumentos válidos, no se rechaza] $\checkmark$} \\
\bottomrule
\end{tabular}%
}
\begin{flushleft}
\footnotesize \textit{Nota.} Variable endógena: EMBI+. Instrumentos excluidos: VIX (serie real desde 1990) y BOVESPA\_VOL (volatilidad realizada trimestral del índice Bovespa, serie real desde 1994), en reemplazo del diferencial construido sobre el ETF EMB, que no cubre la ventana única por cotizar recién desde diciembre de 2007. Errores estándar robustos a heterocedasticidad y autocorrelación (HAC, kernel Newey-West). $^{*}p<0.10$, $^{**}p<0.05$, $^{***}p<0.01$. Elaboración propia (\texttt{codigo/modelos/fase4\_variables\_instrumentales.py}).
\end{flushleft}
\end{table}

El estadístico F de la regresión de primera etapa mide la fuerza de la correlación entre los instrumentos excluidos (VIX, BOVESPA\_VOL) y la variable endógena (EMBI+); alcanza $15.87$ ($p<0.001$), por encima del umbral de referencia de 10 de \textcite{staiger1997} y con mayor margen que el $13.00$ que arrojaba la especificación anterior sobre el ETF EMB: el instrumento nuevo es, en el sentido de la relevancia de primera etapa, más fuerte, no más débil, pese a cambiar de familia de activos (renta variable regional en lugar de renta fija soberana).

El test de sobreidentificación de Sargan no rechaza la validez conjunta de los instrumentos, con un margen aún mayor que antes ($0.006$, $p=0.939$ frente a $0.608$, $p=0.436$): VIX y BOVESPA\_VOL no violan la restricción de exclusión. El test de endogeneidad de Wu-Hausman, en cambio, se debilita de manera sustantiva: ya no rechaza la exogeneidad del EMBI+ ($2.17$, $p=0.144$, frente a $7.32$, $p=0.009$ con el instrumento anterior). Esta combinación, instrumentos más fuertes y más cómodamente válidos, pero endogeneidad formalmente menos sostenida, admite una lectura cautelosa: BOVESPA\_VOL, al ser un instrumento de mercado de renta variable y no de deuda soberana, podría estar captando un componente del apetito de riesgo regional menos ligado específicamente al canal de causalidad inversa (peor resultado fiscal argentino $\to$ mayor EMBI+) que el instrumento de renta fija original, sin que esto invalide el diseño: el Sargan sigue sin rechazar la validez conjunta.

Conviene precisar, además, el alcance informativo del test de Sargan en este contexto. Como señalan \citet{bound1995}, el test de sobreidentificación carece de potencia para detectar violaciones de la restricción de exclusión cuando los instrumentos comparten el mismo canal de transmisión potencial hacia la variable dependiente (en este caso, el apetito de riesgo global que tanto el VIX como el Bovespa capturan). El resultado no significativo del Sargan es condición necesaria pero no suficiente para garantizar la restricción de exclusión: constituye evidencia favorable a la validez de los instrumentos en la medida en que ambos difieran en su canal de transmisión (el VIX mide volatilidad implícita de acciones globales; BOVESPA\_VOL mide volatilidad realizada de la renta variable brasileña), diferencia que el test explota para su identificación, sin eliminar la posibilidad de un canal compartido no detectado. Esta calificación no revierte el diseño instrumental adoptado, sino que delimita el alcance de la evidencia que el Sargan provee.

El coeficiente puntual del EMBI+ instrumentado es positivo ($0.0012$) y alcanza significatividad al 5\% ($p=0.024$), más nítida que la significatividad marginal al 10\% de la especificación anterior ($p=0.076$): un incremento del riesgo soberano instrumentado se asocia con un ajuste positivo, aunque económicamente pequeño, del esfuerzo fiscal. Esta tesis no fuerza una única lectura de conjunto entre un coeficiente más significativo y una prueba de endogeneidad más débil: son dos caras del mismo cambio de instrumento, y ambas se trasladan sin editar al Capítulo \ref{ch:discusion}, que integra esta señal con el resto de la evidencia del capítulo, en particular la ausencia de reacción significativa en DOLS, VECM y Hansen en niveles, antes de extraer una conclusión sobre $H_{1b}$.

El diagnóstico gráfico de primera etapa (EMBI+ real vs. ajustado) y el contraste de causalidad de Granger bivariado entre EMBI+ y $\Delta pb_t$, ambos parte de esta sección en una especificación anterior, no se recalcularon todavía con el instrumento y la ventana nuevos: quedan pendientes de esa actualización y no se reproducen aquí para no mostrar una figura inconsistente con la Tabla \ref{tab:iv2sls}.

\section{Modelo de Umbral de Hansen}
\label{sec:hansen}

El ejercicio más exigente del protocolo econométrico consistió en testear si la Función de Reacción Fiscal exhibe un quiebre de régimen endógeno condicionado por el nivel de riesgo soberano, sin presuponer si ese quiebre toma la forma de una atenuación (fatiga fiscal) o de una reacción ausente en la normalidad y gatillada solo por un nivel extremo de riesgo, siguiendo la metodología de regresión de umbral para series de tiempo de \citet{hansen1999,hansen2000}. El algoritmo realiza una búsqueda secuencial del umbral $\tau$ sobre los percentiles 15 a 85 de la distribución del EMBI+, seleccionando aquel valor que minimiza la suma de cuadrados de los residuos del modelo particionado.

Sobre la ventana única (1996--2025, $n=119$), la búsqueda identificó un umbral óptimo en $\tau^*=2462.1$ puntos básicos de EMBI+, del mismo orden que los $\approx2080$--$2104$ pb que arrojan la especificación con $\Delta pb_t$ (Sección \ref{sec:robustez_hansen_dpb}) y el TVECM (Sección \ref{sec:tvecm_resultados}). La Tabla \ref{tab:hansen} reporta la estimación del modelo particionado en dicho punto.

\begin{table}[H]
\centering
\caption{\label{tab:hansen} Modelo de Umbral de Hansen, $\tau^*=2462.1$ pb de EMBI+, Ventana Única ($n=119$, errores HAC).}
\begin{tabular}{lrrrr}
\toprule
\textbf{Variable} & \textbf{Coef.} & \textbf{Error Est.} & \textbf{$z$} & \textbf{$p$-valor} \\
\midrule
Constante & -0.3830 & 0.570 & -0.672 & 0.502 \\
Brecha del Producto ($\tilde{y}_t$) & 0.0375 & 0.030 & 1.237 & 0.216 \\
$d_{t-1} \cdot \mathbb{I}(\text{EMBI}_t \le 2462.1)$ (Régimen Normal) & -0.0044 & 0.010 & -0.433 & 0.665 \\
$d_{t-1} \cdot \mathbb{I}(\text{EMBI}_t > 2462.1)$ (Régimen de Estrés) & 0.0201 & 0.007 & 2.759 & 0.006$^{***}$ \\
\midrule
\bottomrule
\end{tabular}
\begin{flushleft}
\footnotesize \textit{Nota.} $R^2=0.248$. $R^2$ ajustado$=0.228$. $F=6.287$ ($p<0.001$). Durbin-Watson$=0.364$. $^{***}p<0.01$. Elaboración propia (\texttt{codigo/modelos/fase5\_umbral\_hansen.py}).
\end{flushleft}
\end{table}

El patrón de coeficientes de la ventana única no es el de fatiga fiscal: no es una reacción que se atenúa bajo estrés, es una reacción que directamente no existe en la normalidad y aparece, fuerte y significativa, solo por encima del umbral extremo. En el régimen ``normal'' (EMBI$\le2462$ pb, que cubre 96 de 119 trimestres) el coeficiente es negativo y no significativo ($-0.0044$, $p=0.665$): no hay evidencia de que el resultado primario reaccione a la deuda heredada mientras el riesgo soberano se mantiene por debajo del umbral. En el régimen de ``estrés'' (23 trimestres), en cambio, el coeficiente es positivo y estadísticamente significativo al 1\% ($0.0201$, $p=0.006$): una vez que el riesgo soberano cruza ese nivel extremo, el resultado primario sí reacciona, y con fuerza, a la deuda heredada. Es el patrón inverso al que predice la literatura de fatiga fiscal \citep{ghosh2013}: en vez de una disciplina fiscal que se erosiona bajo presión, aparece una disciplina fiscal que no existe en la normalidad y que la crisis misma gatilla.

El resultado del contraste de significatividad del propio umbral, implementado mediante \textit{bootstrap} del estadístico Sup-LM con 1.000 réplicas, es contundente: el estadístico F observado para la partición óptima es de $0.1857$, con un $p$-valor \textit{bootstrap} menor a $0.001$, que rechaza categóricamente la hipótesis nula de linealidad. A diferencia de la partición extrema que arrojaba la ventana de 88 observaciones (14 contra 74 trimestres, con el régimen minoritario aportando muy poca información), la partición de la ventana única es menos desigual (96 contra 23) y ambos coeficientes de régimen se estiman con mayor precisión, lo que explica que aquí sí resulte significativo el coeficiente del régimen de estrés.

En la especificación en niveles sobre la ventana única, la evidencia de esta tesis no ofrece una señal de fatiga fiscal, ofrece una señal de régimen de reacción fiscal extremo: ausencia de reacción en la normalidad, reacción significativa solo bajo estrés agudo. La Sección \ref{sec:robustez_hansen_dpb} evalúa el mismo fenómeno bajo una especificación alternativa que satisface el supuesto de estacionariedad del test, y sobre la ventana de 88 observaciones (pendiente de reestimarse sobre la ventana única), con un patrón de atenuación más cercano a la fatiga fiscal clásica: el Capítulo \ref{ch:discusion} desarrolla las implicancias de que las dos especificaciones, sobre ventanas distintas, apunten a formas cualitativamente distintas de no linealidad pese a coincidir en la ubicación del umbral.

\textit{Nota metodológica sobre el supuesto de estacionariedad del test de umbral.} El contraste Sup-LM y su \textit{bootstrap} son asintóticamente válidos bajo el supuesto de estacionariedad estricta de todos los regresores, incluida la variable dependiente \citep[Assumption~A, p.~577]{hansen2000}. Sin embargo, el diagnóstico de la Sección \ref{sec:dfgls} clasifica $pb_t$ como I(1), de modo que dicho supuesto no se cumple en sentido estricto en la especificación en niveles reportada arriba. \citet{caner2004} extienden el marco de umbral al caso de procesos cointegrados. El resultado reportado permanece como el diagnóstico central de esta sección; la Sección \ref{sec:robustez_hansen_dpb} reporta, como verificación adicional, la misma especificación con la variable dependiente en primera diferencia (I(0)).

\subsection{Robustez: Umbral de Hansen con la Variable Dependiente en Primera Diferencia}
\label{sec:robustez_hansen_dpb}

La nota metodológica anterior exige reestimar el modelo de umbral bajo una especificación cuya variable dependiente satisfaga el supuesto de estacionariedad estricta de \citet{hansen2000}. Dado que $pb_t$ es I(1) mientras que su primera diferencia $\Delta pb_t$ es I(0), se reestimó exactamente el mismo procedimiento reemplazando $pb_t$ por $\Delta pb_t$ como variable dependiente. Esta especificación no se recalculó todavía sobre la ventana única: los resultados que siguen corresponden a la ventana de 88 observaciones, pendientes de esa actualización.

\begin{table}[H]
\centering
\caption{\label{tab:hansen_dpb} Modelo de Umbral de Hansen con $\Delta pb_t$ (I(0)) como Dependiente ($n=87$, errores HAC).}
\begin{tabular}{lrrrr}
\toprule
\textbf{Variable} & \textbf{Coef.} & \textbf{Error Est.} & \textbf{$z$} & \textbf{$p$-valor} \\
\midrule
Constante & -1.0673 & 0.290 & -3.682 & 0.000 \\
Brecha del Producto ($\tilde{y}_t$) & -0.0019 & 0.013 & -0.151 & 0.880 \\
$d_{t-1} \cdot \mathbb{I}(\text{EMBI}_t \le 2083)$ (Régimen Normal) & 0.0184 & 0.005 & 3.412 & 0.001 \\
$d_{t-1} \cdot \mathbb{I}(\text{EMBI}_t > 2083)$ (Régimen de Estrés) & 0.0077 & 0.004 & 2.166 & 0.030 \\
\midrule
\bottomrule
\end{tabular}
\begin{flushleft}
\footnotesize \textit{Nota.} Umbral óptimo $\tau^*=2083$ pb de EMBI+ (frente a 449.4 pb en la especificación en niveles; partición 74 contra 14 observaciones, inversa a la de la especificación en niveles). \textit{Bootstrap} Sup-LM (1.000 réplicas): $F=0.2326$, $p<0.001$. Elaboración propia.
\end{flushleft}
\end{table}

El resultado se revierte de manera sustancial respecto de la especificación en niveles. El \textit{bootstrap} Sup-LM rechaza la linealidad al 1\% ($p<0.001$, frente a $p=0.076$ en niveles), y ambos coeficientes de régimen resultan individualmente significativos: $0.0184$ ($p=0.001$) en el régimen normal y $0.0077$ ($p=0.030$) en el régimen de estrés, con el patrón exacto que predeciría la fatiga fiscal (reacción atenuada, no ausente, bajo estrés financiero). A diferencia de la especificación en niveles (Sección \ref{sec:hansen}), aquí la partición no es extrema en el mismo sentido: el umbral de $2083$ pb ubica al régimen ``normal'' como el mayoritario (74 de 88 trimestres) y al de ``estrés'' como el minoritario (14), la partición inversa a la hallada en niveles, pero igualmente desigual en magnitud.

Esta especificación exige, sin embargo, precisar su alcance antes de extraer una conclusión sustantiva. El objeto que estima no es idéntico al de la Tabla \ref{tab:hansen}: mientras aquella modela el \textit{nivel} del resultado primario en función del \textit{nivel} de la deuda rezagada por régimen (la forma canónica de \citet{ghosh2013}), esta especificación modela la \textit{variación} trimestral del resultado primario frente al mismo regresor, un objeto más cercano a un mecanismo de ajuste de corto plazo que a la regla de reacción de nivel de la literatura de referencia.

Con esa salvedad explícita, el resultado es sustantivo: bajo la especificación que satisface el supuesto de estacionariedad exigido por el test, la evidencia de un umbral, y de un patrón de atenuación coherente con fatiga fiscal, es estadísticamente significativa. Las tres especificaciones de umbral de esta tesis (niveles sobre la ventana única, $\Delta pb_t$ sobre la ventana de 88 observaciones, y el TVECM sobre la ventana única) coinciden en rechazar la linealidad y en ubicar el umbral en un rango estrecho ($\tau^*\approx2080$--$2462$ pb), pero no coinciden en la forma que toma la no linealidad: niveles encuentra una reacción ausente en la normalidad y gatillada solo por estrés extremo, mientras que $\Delta pb_t$ encuentra una reacción presente en ambos regímenes que se atenúa, sin desaparecer, bajo estrés. Esta divergencia en la forma funcional, pese a la convergencia en la ubicación del umbral, no debe sobreinterpretarse como dos confirmaciones independientes de un mismo fenómeno: son objetos distintos (nivel de $pb_t$ frente a su variación trimestral) sobre ventanas distintas (única frente a 88 observaciones), y la comparación exige la cautela adicional de que $\Delta pb_t$ está pendiente de reestimarse sobre la ventana única. El Capítulo \ref{ch:discusion} retoma esta tensión entre especificaciones para calificar el diagnóstico de $H_{1a}$ del Capítulo \ref{ch:conclusiones}.

\section{Sensibilidad Temporal: Reestimación por Subperíodos}
\label{sec:sensibilidad_temporal}

Como ejercicio de robustez adicional, se reestimó la relación entre el Resultado Primario, la deuda rezagada y la Brecha del Producto en dos submuestras de tamaño similar, 2004T1--2014T4 ($n=43$) y 2015T1--2025T4 ($n=44$), evaluando si el coeficiente de reacción fiscal es sensible a la ventana temporal considerada. Se reportan tres especificaciones por subperíodo: la estática de Mínimos Cuadrados Ordinarios con errores HAC, y dos versiones DOLS con ampliación dinámica reducida ($m=1$ y $m=2$ adelantos y rezagos, frente a los $m=4$ del DOLS de muestra completa de la Tabla \ref{tab:dols}). La reducción del orden de ampliación es necesaria para preservar grados de libertad en submuestras de aproximadamente 40 observaciones.

\begin{table}[H]
\centering
\caption{\label{tab:sensibilidad_temporal} Sensibilidad Temporal del Coeficiente de Reacción Fiscal ($\rho$) por Subperíodo: Especificación Estática y DOLS.}
\begin{tabular}{llrrrr}
\toprule
\textbf{Subperíodo} & \textbf{Especificación} & \textbf{$n$} & \textbf{$\rho$ ($d_{t-1}$)} & \textbf{Error Est. (HAC)} & \textbf{$p$-valor} \\
\midrule
\multirow{3}{*}{2004T1--2014T4} & MCO estático & 43 & 0.0373 & 0.0118 & 0.0015 \\
 & DOLS ($m=1$) & 40 & 0.0384 & 0.0207 & 0.0632 \\
 & DOLS ($m=2$) & 38 & 0.0604 & 0.0409 & 0.1398 \\
\midrule
\multirow{3}{*}{2015T1--2025T4} & MCO estático & 44 & 0.0530 & 0.0176 & 0.0026 \\
 & DOLS ($m=1$) & 41 & 0.0364 & 0.0175 & 0.0377 \\
 & DOLS ($m=2$) & 39 & 0.0357 & 0.0183 & 0.0508 \\
\midrule
\textit{Muestra completa} & \textit{DOLS ($m=4$)} & \textit{78} & \textit{-0.0071} & \textit{0.0123} & \textit{0.5640} \\
\bottomrule
\end{tabular}
\begin{flushleft}
\footnotesize \textit{Nota.} Especificación estática: $pb_t = \alpha + \rho\, d_{t-1} + \gamma\, \tilde{y}_t + \varepsilon_t$. Especificación DOLS: la anterior más $\sum_{j=-m}^{m}\phi_j \Delta d_{t-j}$. Errores HAC (Newey-West) en todos los casos. La fila final reproduce, para comparación, la estimación de muestra completa de la Tabla \ref{tab:dols}. Elaboración propia.
\end{flushleft}
\end{table}

De ello se siguen dos conclusiones centrales: primero, parte de la magnitud del coeficiente estático responde al sesgo de simultaneidad de corto plazo que la ampliación de Stock-Watson corrige (en el segundo subperíodo cae de 0.0530 a 0.0364); segundo, el signo positivo no es un mero artefacto estático. La discrepancia con el DOLS de muestra completa ($\rho=-0.0071$) responde a la inestabilidad paramétrica de la relación a lo largo del tiempo: dos submuestras con coeficiente positivo coexisten con un coeficiente agregado nulo si el parámetro no es constante, en consonancia con los quiebres estructurales múltiples de Bai-Perron documentados en la Sección \ref{sec:quiebres_estructurales} del Capítulo \ref{ch:descriptivo}.

El resultado nulo de muestra completa debe entenderse como la ausencia de una regla de reacción fiscal \textit{estable a lo largo de 2004--2025}. La magnitud de la respuesta por subperíodo sigue siendo económicamente pequeña (entre 0.036 y 0.060 puntos del PIB por punto de deuda), insuficiente para estabilizar la dinámica de endeudamiento.

\section{Resultados del VAR Estructural Restringido (SVAR)}
\label{sec:svar_resultados}

La estimación del Vector Autoregresivo Estructural Restringido (SVAR) con reglas macroeconómicas de Blanchard-Perotti ($\eta_{pb} = 0.25$, $a_{23}=a_{24}=0$) y Favero-Giavazzi sobre el vector de cinco variables $\mathbf{Y}_t = (\tilde{y}_t, pb_t, \text{EMBI}_t, \text{TCRM}_t, d_t)'$ permite desentrañar la transmisión contemporánea y dinámica de las perturbaciones macro-fiscales. Se estima sobre la ventana ampliada de treinta años (1996--2025, $n=120$), no sobre la ventana original: es la especificación de referencia del resto del protocolo (Sección \ref{sec:vecm} del Capítulo \ref{ch:metodologia}), y su resultado cambia de manera sustantiva respecto de estimarla solo sobre 2004--2025, como se explicita más abajo. La Tabla \ref{tab:svar_fevd} reporta la Descomposición de Varianza del Error de Pronóstico (FEVD) para horizontes de 1 a 20 trimestres.

\begin{table}[H]
\centering
\caption{\label{tab:svar_fevd} Descomposición de Varianza del Error de Pronóstico (FEVD) del SVAR Restringido, Ventana Ampliada 1996--2025 (\%).}
\resizebox{\textwidth}{!}{%
\begin{tabular}{llrrrrr}
\toprule
\textbf{Horizonte} & \textbf{Variable Explicada} & \textbf{Shock Brecha PIB} & \textbf{Shock Superávit} & \textbf{Shock EMBI+} & \textbf{Shock TCRM} & \textbf{Shock Deuda Propia} \\
\midrule
\multirow{2}{*}{1 Trimestre} & Deuda Pública ($d_t$) & 20.89 & 23.94 & 3.23 & 8.91 & 43.03 \\
 & Riesgo País ($\text{EMBI}_t$) & 7.81 & 11.30 & 80.89 & 0.00 & 0.00 \\
\midrule
\multirow{2}{*}{4 Trimestres} & Deuda Pública ($d_t$) & 12.59 & 18.97 & 10.13 & 5.39 & 52.92 \\
 & Riesgo País ($\text{EMBI}_t$) & 7.02 & 7.34 & 70.64 & 0.42 & 14.58 \\
\midrule
\multirow{2}{*}{8 Trimestres} & Deuda Pública ($d_t$) & 8.14 & 14.90 & 12.59 & 3.45 & 60.91 \\
 & Riesgo País ($\text{EMBI}_t$) & 8.44 & 7.81 & 53.10 & 0.72 & 29.93 \\
\midrule
\multirow{2}{*}{20 Trimestres} & Deuda Pública ($d_t$) & 7.28 & 12.70 & 11.36 & 14.44 & 54.21 \\
 & Riesgo País ($\text{EMBI}_t$) & 14.92 & 16.16 & 34.09 & 9.80 & 25.04 \\
\bottomrule
\end{tabular}%
}
\begin{flushleft}
\footnotesize \textit{Nota.} Identificación estructural por restricciones macroeconómicas de Blanchard-Perotti ($\eta_{pb}=0.25$, rigidez intra-trimestral en superávit) e identidad de acumulación de Favero-Giavazzi. Valores extraídos directamente de \texttt{resultados/tablas/fase19\_fevd.csv} (ventana 1996--2025; verificados sin cambio sustantivo contra la ventana real 1997--2025 sin trimestres aproximados, \texttt{fase23b\_fevd\_1997\_2025.csv}).
\end{flushleft}
\end{table}

Extender la ventana de 26 a 30 años revierte la lectura que esta tesis sostenía sobre la ventana original (2004--2025): allí, la descomposición de varianza a 20 trimestres atribuía la mayor parte de la variabilidad de la deuda a los shocks fiscal y de actividad (47{,}5\% y 28{,}6\% respectivamente, con el canal cambiario y de riesgo soberano acotado al 11{,}5\% conjunto). Sobre la ventana ampliada, el patrón se invierte: la inercia propia de la deuda pasa a explicar el 54{,}2\% de su propia varianza a 20 trimestres (12{,}4\% en la ventana original), el canal cambiario y de riesgo soberano conjunto (25{,}8\%) supera ahora al fiscal y de actividad conjunto (20{,}0\%), y ningún shock individual domina de forma aislada como lo hacía el superávit primario antes. Esta reversión no es un artefacto de la única aproximación de datos que introduce la ventana ampliada, la extrapolación del TCRM de 1996 (Sección \ref{sec:empalme} del Capítulo \ref{ch:metodologia}): reestimada sobre la ventana real sin esos cuatro trimestres (1997--2025, $n=116$), la inercia propia sigue explicando el 53{,}2\% y el canal cambiario-riesgo el 27{,}5\%, prácticamente idéntico. Es, en cambio, coherente con el resultado de la Sección \ref{sec:resultados_vecm}: la ventana ampliada muestra menos evidencia de una relación de cointegración estable que la original, y un sistema con menor capacidad de auto-corrección hacia un equilibrio de largo plazo es, por construcción, un sistema donde la inercia de cada variable explica más de su propia varianza futura. El Capítulo \ref{ch:discusion} retoma esta reversión para evaluar si treinta años de evidencia, en lugar de veintiséis, sostienen mejor o peor la lectura sustantiva de esta tesis sobre el canal dominante de ajuste de la deuda.

\begin{figure}[H]
\centering
\includegraphics[width=0.92\textwidth]{figura_svar_irf.png}
\caption{Funciones de Impulso-Respuesta Estructurales (SVAR Restringido) con Bandas Bootstrap al 95\%.}
\label{fig:svar_irf}
\end{figure}

\begin{figure}[H]
\centering
\includegraphics[width=0.88\textwidth]{figura_svar_fevd.png}
\caption{Descomposición de Varianza del Error de Pronóstico (FEVD) de Deuda/PIB y Riesgo País.}
\label{fig:svar_fevd}
\end{figure}

\section{Estimación del Modelo Threshold VECM (TVECM)}
\label{sec:tvecm_resultados}

Para verificar si la velocidad de corrección fiscal conmuta entre regímenes de estrés soberano en un marco multivariado, se estimó el modelo Threshold VECM de \citet{hansen2002} sobre las cuatro variables del sistema de cointegración ($d_t, pb_t, \text{EMBI}_t, \text{TCRM}_t$), sobre la misma ventana ampliada de treinta años (1996--2025) que sostiene el VECM y el SVAR restringido.

\begin{table}[H]
\centering
\caption{\label{tab:tvecm} Estimación del Modelo Threshold VECM (Hansen y Seo, 2002), Ventana Ampliada 1996--2025.}
\begin{tabular}{lrr}
\toprule
\textbf{Parámetro / Estadístico} & \textbf{Régimen 1: Normalidad} & \textbf{Régimen 2: Fatiga / Estrés} \\
\midrule
Umbral Óptimo de Riesgo ($\tau^*$) & \multicolumn{2}{c}{$\mathbf{2104.22\text{ pb}}$ (Soporte empírico: 15\%--85\%)} \\
Observaciones Efectivas por Régimen & 96 trimestres ($81.4\%$) & 22 trimestres ($18.6\%$) \\
Velocidad de Ajuste Fiscal ($\alpha_{pb}$) & $\mathbf{+0.0011}$ & $\mathbf{-0.0074}$ \\
Estadístico Sup-LM de Linealidad & \multicolumn{2}{c}{$\text{Sup-LM} = 59.78 \quad (p\text{-valor bootstrap} = 0.632)$} \\
\bottomrule
\end{tabular}
\begin{flushleft}
\footnotesize \textit{Nota.} Vector de cointegración estimado por Johansen. Datos provenientes de \texttt{resultados/tablas/fase23\_tvecm\_resumen.csv}. El umbral estimado ($\tau^*=2104$ pb) es prácticamente idéntico al de la ventana original 2004--2025 ($\tau^*=2081$ pb) y al del umbral de Hansen sobre $\Delta pb_t$ (Sección \ref{sec:robustez_hansen_dpb}, $\tau^*=2083$ pb): las tres especificaciones, sobre ventanas y variables dependientes distintas, coinciden en un umbral de riesgo soberano de entre 2080 y 2105 puntos básicos.
\end{flushleft}
\end{table}

La reversión de signo entre regímenes ($\alpha_{pb}$ pasa de $+0{,}0011$ bajo estrés bajo a $-0{,}0074$ bajo estrés alto) se replica con la misma dirección que en la ventana original, aunque con magnitudes más pequeñas en el régimen normal; el estadístico Sup-LM sigue sin alcanzar significatividad ($p=0{,}632$, frente a $p=0{,}625$ en la ventana original), de modo que la lectura de esta sección no cambia: hay un patrón de umbral notablemente estable en su localización entre especificaciones, pero no una prueba estadística formal de que la conmutación de régimen sea significativa bajo este modelo multivariado en particular.

\section{Simulación Estocástica del DSA con Covarianza del SVAR Restringido}
\label{sec:dsa_resultados}

La especificación de referencia del Análisis de Sostenibilidad de la Deuda estocástico (Sección \ref{sec:dsa_metodologia} del Capítulo \ref{ch:metodologia}) sustituye el proceso CIR y el modelo DCC-GARCH de una especificación anterior por una matriz de correlación derivada directamente de la covarianza de residuos estructurales del SVAR restringido de la Sección \ref{sec:svar_resultados}, estimado sobre la ventana ampliada de treinta años. Bajo esta especificación, la probabilidad de que la deuda consolidada supere el 100\% del PIB hacia 2035 es $\mathbf{26{,}1\%}$, con una mediana de $82{,}8\%$ del PIB en el escenario de Referencia determinista. La Figura \ref{fig:fan_chart_final} (Capítulo \ref{ch:discusion}) presenta el gráfico de abanico correspondiente. Es un resultado del mismo orden de magnitud que una especificación anterior calibrada a mano ($31{,}2\%$, con $29{,}8\%$ bajo un proceso CIR del riesgo soberano y $32{,}5\%$ bajo DCC-GARCH), documentada aparte en un memo de trabajo para el director de tesis y no en este documento, pero no idéntico: la matriz de correlación fundada en el SVAR estimado, a diferencia de la calibrada a mano, incorpora directamente el hallazgo de la Sección \ref{sec:svar_resultados} de que el riesgo soberano y el tipo de cambio real están, sobre esta ventana, menos correlacionados entre sí de lo que asumía la calibración original, lo que angosta moderadamente la cola derecha de la distribución simulada.

\section{Extensión de Frontera: Reconstrucción del Spread Soberano en 42 Años de Democracia (1983--2025)}
\label{sec:spread_historico_1983_2025}

Para evaluar si la dinámica de riesgo crediticio soberano y los umbrales de fatiga fiscal constituyen fenómenos circunscritos a las últimas dos décadas o si representan regularidades empíricas invariantes de la economía política argentina, se procedió a reconstruir y empalmar una serie trimestral homogénea de costo de financiamiento soberano a lo largo de los 42 años de democracia ininterrumpida ($1983\text{T4}$--$2025\text{T4}$, $n = 169$ trimestres).

\subsection{Metodología de Reconstrucción por Eras Financieras}

Ante la inexistencia del índice EMBI con anterioridad a 1993, la serie histórica se estructuró a partir de los activos soberanos en moneda extranjera que fijaron el costo marginal de endeudamiento en cada subperíodo:
\begin{enumerate}
    \item \textbf{Era Bonex (1983T4--1992T4)}: Se utilizaron los rendimientos secundarios en dólares de los Bonos Externos de la República Argentina (Series 1982, 1984, 1987 y 1989) cotizados en el Mercado Abierto Electrónico y la Bolsa de Comercio de Buenos Aires, deduciendo la tasa del Bono del Tesoro de EE.UU. a 10 años ($\text{Spread}_t = \text{TIR}_{\text{Bonex}, t} - r_{US, 10Y, t}$), conforme a la metodología canónica de \citet{neumeyer2005} y \citet{kehoe2021}.
    \item \textbf{Era Brady (1993T1--1997T4)}: Se adoptó el JP Morgan EMBI Argentina \textit{Stripped Spread} sobre bonos Brady Par, Discount y FRB \citep{uribe2006}, descontando el valor presente de los bonos cupón cero del Tesoro estadounidense que colateralizaban el capital.
    \item \textbf{Era EMBI+ y Global (1998T1--2025T4)}: Se integró la serie diaria oficial de JP Morgan EMBI+ y EMBI Global Diversified Argentina.
\end{enumerate}

Para evitar discontinuidades espurias en los puntos de unión, se aplicó una corrección proporcional de sesgo en solapamiento (\textit{Overlap Bias Correction}), vinculando la escala hacia el benchmark moderno del EMBI+:
\begin{equation}
k_1 = \frac{\overline{\text{EMBI}}_{\text{Brady, 1993T1}}}{\overline{\text{Spread}}_{\text{Bonex, 1992T4}}} = \frac{850.0}{760.0} = 1.1184, \quad k_2 = \frac{\overline{\text{EMBI+}}_{\text{1998T1}}}{\overline{\text{EMBI}}_{\text{Brady, 1997T4}}} = \frac{540.0}{620.0} = 0.8710
\end{equation}

\begin{figure}[H]
\centering
\includegraphics[width=\textwidth]{figura_spread_historico_1983_2025.png}
\caption{Evolución Histórica del Spread Soberano Argentino en 42 Años de Democracia (1983--2025).}
\label{fig:spread_historico}
\end{figure}

\subsection{Quiebres Estructurales Múltiples de Bai-Perron en el Ciclo Histórico}

La aplicación del algoritmo de programación dinámica de \citet{bai2003} sobre los 169 trimestres de democracia identifica cinco quiebres mayores y seis regímenes de financiamiento soberano con diferencias de media altamente significativas (Tabla \ref{tab:bai_perron_historico}).

\begin{table}[H]
\centering
\caption{\label{tab:bai_perron_historico} Quiebres Estructurales Múltiples de Bai-Perron en el Spread Soberano (1983--2025).}
\resizebox{\textwidth}{!}{%
\begin{tabular}{lllrrrl}
\toprule
\textbf{Régimen} & \textbf{Fecha Inicio} & \textbf{Fecha Fin} & \textbf{Trimestres} & \textbf{Media (pb)} & \textbf{Desvío (pb)} & \textbf{Hito Histórico / Ruptura} \\
\midrule
Régimen 1 & 1983T4 & 1988T4 & 21 & 797.84 & 359.44 & Transición democrática y colapso del Plan Austral \\
Régimen 2 & 1989T1 & 1990T4 & 8 & 3573.76 & 1255.84 & Hiperinflaciones de 1989/1990 y Plan Bonex \\
Régimen 3 & 1991T1 & 2001T4 & 44 & 831.66 & 427.31 & Ley de Convertibilidad y Plan Brady \\
Régimen 4 & 2002T1 & 2005T2 & 14 & 5431.33 & 638.09 & Default de la deuda externa y crisis de 2001--2002 \\
Régimen 5 & 2005T3 & 2019T2 & 56 & 683.54 & 326.60 & Canjes 2005/2010, superávits gemelos y desendeudamiento \\
Régimen 6 & 2019T3 & 2025T4 & 26 & 1732.05 & 632.45 & Shock PASO 2019, Canje 2020 y estabilización post-2023 \\
\bottomrule
\end{tabular}%
}
\begin{flushleft}
\footnotesize \textit{Nota.} Estimación por minimización de pérdida L2 mediante programación dinámica sobre 169 observaciones trimestrales. Datos de \texttt{resultados/tablas/fase22\_bai\_perron\_1983\_2025.csv}.
\end{flushleft}
\end{table}

\subsection{Calibración del Proceso CIR y Métricas por Régimen Presidencial}

La calibración por Máxima Verosimilitud exacta del proceso CIR a lo largo de las cuatro décadas (Tabla \ref{tab:cir_historico}) demuestra que la velocidad de reversión a la media de largo plazo es $\kappa = 0.4630$ (vida media de $1.50$ años), con un nivel incondicional de equilibrio de $\theta = 1456.98$ pb. La condición de Feller se satisface de forma holgada ($2\kappa\theta = 1349.08 > \sigma^2 = 613.96 \implies \text{Ratio} = 2.1974 > 1.0$), ratificando la robustez del modelado en tiempo continuo.

\begin{table}[H]
\centering
\caption{\label{tab:cir_historico} Calibración MLE del Proceso CIR en la Muestra Ultra-Larga (1983--2025).}
\begin{tabular}{lrr}
\toprule
\textbf{Parámetro / Métrica} & \textbf{Valor Estimado (42 Años, $n=169$)} & \textbf{Interpretación Económica} \\
\midrule
Velocidad de Reversión ($\kappa$) & $\mathbf{0.4630}$ & Reversión media hacia equilibrio \\
Nivel Incondicional ($\theta$) & $\mathbf{1456.98\text{ pb}}$ & Costo medio estructural de endeudamiento \\
Volatilidad de Difusión ($\sigma$) & $\mathbf{24.7781}$ & Dispersión de innovaciones estocásticas \\
Vida Media de Reversión ($t_{1/2}$) & $\mathbf{1.50\text{ años}}$ ($18$ meses) & Duración típica de los ciclos de spread \\
Condición de Feller ($2\kappa\theta > \sigma^2$) & $1349.08 > 613.96 \implies \mathbf{Cumple}$ & No negatividad estricta garantizada \\
Ratio de Feller ($2\kappa\theta / \sigma^2$) & $\mathbf{2.1974 > 1.0}$ & Frontera en cero inaccesible ($\mathbb{P}=1$) \\
Criterio AIC (CIR vs. AR(1)) & $\mathbf{2466.56} < 2612.53$ & $\Delta \text{AIC} = -145.97 \implies$ Prefiere CIR \\
\bottomrule
\end{tabular}
\begin{flushleft}
\footnotesize \textit{Nota.} Estimación exacta sobre la función de densidad Chi-cuadrado no central. Datos extraídos de \texttt{resultados/tablas/fase22\_cir\_calibracion\_1983\_2025.csv}.
\end{flushleft}
\end{table}

\begin{table}[H]
\centering
\caption{\label{tab:spread_por_regimen} Dinámica del Spread Soberano por Régimen Presidencial (1983--2025).}
\resizebox{\textwidth}{!}{%
\begin{tabular}{llrrrrr}
\toprule
\textbf{Administración Presidencial} & \textbf{Período} & \textbf{Trimestres} & \textbf{Media (pb)} & \textbf{Mediana (pb)} & \textbf{Mínimo (pb)} & \textbf{Máximo (pb)} \\
\midrule
Raúl Alfonsín & 1983T4--1989T2 & 23 & 1035.52 & 701.36 & 409.13 & 4675.72 \\
Carlos Menem I & 1989T3--1995T1 & 23 & 1577.73 & 837.73 & 566.13 & 5455.01 \\
Carlos Menem II & 1995T2--1999T3 & 18 & 688.63 & 663.90 & 357.10 & 1350.00 \\
Fernando de la Rúa & 1999T4--2001T3 & 8 & 811.26 & 710.21 & 550.92 & 1468.94 \\
A. Rodríguez Saá / E. Duhalde & 2001T4--2003T1 & 6 & 5354.62 & 5873.34 & 2947.64 & 6658.98 \\
Néstor Kirchner & 2003T2--2007T3 & 18 & 2775.65 & 2644.15 & 210.22 & 5800.85 \\
Cristina Fernández de Kirchner & 2007T4--2015T3 & 32 & 853.48 & 734.62 & 387.78 & 1696.86 \\
Mauricio Macri & 2015T4--2019T3 & 16 & 584.87 & 477.95 & 360.23 & 1493.83 \\
Alberto Fernández & 2019T4--2023T3 & 16 & 2048.85 & 2038.76 & 1391.80 & 3103.00 \\
Javier Milei & 2023T4--2025T4 & 9 & 1195.30 & 870.36 & 684.00 & 2257.75 \\
\bottomrule
\end{tabular}%
}
\begin{flushleft}
\footnotesize \textit{Nota.} Spread empalmado expresado en puntos básicos. Datos de \texttt{resultados/tablas/fase22\_spread\_por\_regimen.csv}.
\end{flushleft}
\end{table}

Este análisis de ultra-largo plazo confirma que el umbral de fatiga fiscal estimado en la tesis ($\tau^* \approx 2080$--$2083$ pb) constituye una cota estructural: los episodios en que el spread superó dicha marca coinciden de manera exacta con los períodos de crisis terminal de la deuda (1989--1990, 2001--2005 y 2020--2023), demostrando que la restricción externa de financiamiento rige la solvencia argentina con independencia del régimen monetario o político vigente.

\section{Interpretación de la Magnitud Económica de los Coeficientes}

Corresponde cuantificar la magnitud económica de los coeficientes puntuales estimados. El coeficiente de reacción fiscal del DOLS de muestra completa ($\rho=-0.0071$) implica que un incremento de 10 puntos porcentuales en la ratio Deuda/PIB rezagada se asocia con una reducción de apenas 0.07 puntos del PIB en el resultado primario: una magnitud minúscula que confirma que la falta de significatividad no enmascara un efecto cuantitativamente relevante. En la estimación IV-2SLS (Sección \ref{sec:iv2sls}), el coeficiente del EMBI+ instrumentado ($0.0017$, marginalmente significativo al 10\%) implica que un incremento de 100 puntos básicos en el riesgo soberano genera una variación de apenas 0.17 puntos del PIB en el superávit primario: incluso bajo la lectura más favorable a la existencia de un canal de reacción vía riesgo soberano, la magnitud sigue siendo pequeña en términos macroeconómicos frente a los desequilibrios fiscales documentados en el Capítulo \ref{ch:descriptivo}.

\section{Síntesis de la Evidencia Empírica}
\label{sec:sintesis_empirica}

Siete estrategias de estimación convergen, todas sobre la ventana muestral única de treinta años (1996--2025, $n=120$), en un diagnóstico predominantemente común sobre la ausencia de una reacción fiscal lineal estable, con dos divergencias sustantivas, la forma del umbral de no linealidad y el canal dominante de transmisión hacia la deuda, que exigen calificación explícita.

Sobre la reacción lineal: el contraste de Johansen no rechaza la ausencia de cointegración entre las cuatro variables del sistema (Sección \ref{sec:johansen_cap6}); DOLS no encuentra una reacción significativa al endeudamiento y, a diferencia de estimaciones sobre ventanas más cortas, el signo puntual ya es el correcto ($\rho=0.0159$, $p=0.218$); el VECM replica el resultado nulo con más claridad aún ($\alpha_{pb}\approx0$, $p=0.996$) y encuentra un canal de ajuste vía riesgo soberano significativo ($\alpha_{EMBI}=6.356$, $p=0.025$); e IV-2SLS, con un instrumento regional reconstruido para cubrir la ventana completa (VIX y la volatilidad del Bovespa, Sección \ref{sec:iv2sls}), encuentra un coeficiente del EMBI+ instrumentado positivo y significativo al 5\% ($0.0012$, $p=0.024$), aunque el test de Wu-Hausman ya no rechaza con fuerza la exogeneidad del EMBI+ ($p=0.144$, frente a $p=0.009$ con el instrumento anterior).

Sobre la forma del umbral de no linealidad: el umbral de Hansen en niveles, la especificación en $\Delta pb_t$ (todavía sobre 88 observaciones, pendiente de reestimarse) y el TVECM coinciden en ubicar un umbral de riesgo soberano en un rango estrecho ($\tau^*\approx2080$--$2462$ pb), pero no coinciden en qué tipo de no linealidad describen: Hansen en niveles encuentra una reacción ausente en la normalidad y gatillada solo por estrés extremo (régimen normal $-0.0044$, no significativo; régimen de estrés $+0.0201$, $p=0.006$), un patrón de régimen extremo y no de fatiga fiscal; Hansen en $\Delta pb_t$ encuentra, en cambio, una reacción presente en ambos regímenes que se atenúa bajo estrés sin desaparecer, más cercana a la fatiga fiscal clásica; y el TVECM encuentra una reversión de signo entre regímenes ($\alpha_{pb}$ pasa de $+0.001$ a $-0.007$) sin alcanzar significatividad en su contraste Sup-LM ($p=0.632$). Tres formas funcionales distintas sobre el mismo umbral, no una convergencia limpia.

Sobre el canal dominante de transmisión: la descomposición de varianza del SVAR restringido atribuye a la inercia propia de la deuda el $54.2\%$ de su propia varianza a 20 trimestres y al canal cambiario-riesgo conjunto el $25.8\%$, superando ahora al fiscal-real conjunto ($20.0\%$); Bai-Perron, corrido por primera vez sobre la ventana completa, detecta tres quiebres estructurales en la deuda SPNF, incluido el de 2001 que una ventana más corta no podía identificar. El DSA estocástico, con una covarianza derivada de esta misma estimación del SVAR en lugar del proceso calibrado a mano de una especificación anterior (documentada aparte, fuera de este documento), sitúa la probabilidad de insolvencia a 2035 en $26.1\%$.

En síntesis: la evidencia de una reacción fiscal lineal, estable y de nivel sigue siendo predominantemente débil o nula, con el signo puntual ya orientado en la dirección correcta aunque sin significatividad (DOLS, VECM); la evidencia de un régimen crítico de riesgo soberano es robusta en su ubicación pero mixta en su forma, entre atenuación y reacción gatillada por crisis extrema, según la especificación; y la lectura sobre el canal dominante de transmisión hacia la deuda es, con treinta años de evidencia, la inercia propia del stock y el riesgo cambiario-soberano, no el desbalance fiscal corriente. El Capítulo \ref{ch:discusion} discute en profundidad estas tres calificaciones antes de evaluar sus implicancias teóricas sobre la solvencia intertemporal.
